% =====================================================================
% Atualizações da proposta do SIGCF para colar no Overleaf.
% Compatível com abnTeX2 (usa apenas itemize, textbf e textit).
% =====================================================================

% ---------------------------------------------------------------------
% 1. Requisitos funcionais — itens a acrescentar
%    RF01 a RF04 permanecem. Os itens abaixo formalizam funcionalidades
%    que já constam da solução descrita e foram implementadas no MVP.
% ---------------------------------------------------------------------

\begin{itemize}
  \item \textbf{RF01 (ampliado) -- Cadastros:} manter também o cadastro de clientes (nome, telefone opcional e endereço), utilizado para sugerir o endereço das ordens de coleta e consultar o histórico de atendimentos por cliente.
  \item \textbf{RF05 -- Rastreabilidade:} registrar cada mudança de status de uma ordem de coleta com data, hora, usuário responsável, status anterior, novo status e motivo, sendo o motivo obrigatório nas transições para Falha e Cancelada.
  \item \textbf{RF06 -- Painel do dia:} apresentar ao gestor a quantidade de coletas do dia por status, atualizada automaticamente, e a lista das coletas com seus responsáveis.
\end{itemize}

% ---------------------------------------------------------------------
% 2. Requisitos não funcionais — substitui a versão anterior
%    (que citava Spring Boot/Java e PostgreSQL)
% ---------------------------------------------------------------------

\begin{itemize}
  \item \textbf{RNF01 -- Interface web responsiva:} a aplicação deve ser uma \textit{Single Page Application} (SPA) desenvolvida em React, com a interface do motorista projetada segundo a abordagem \textit{mobile-first} e instalável como \textit{Progressive Web App} (PWA), dispensando a distribuição por lojas de aplicativos.
  \item \textbf{RNF02 -- API em camadas:} o \textit{back-end} deve expor uma API REST desenvolvida em Node.js com o \textit{framework} NestJS, organizada nas camadas \textit{Controller}, \textit{Service} e \textit{Repository}, com persistência mediada pelo mapeamento objeto-relacional TypeORM.
  \item \textbf{RNF03 -- Banco de dados relacional:} os dados devem ser armazenados no MySQL 8 com o mecanismo InnoDB, que garante transações ACID nas operações logísticas, como a alteração de status de uma ordem e a gravação do registro correspondente no histórico.
  \item \textbf{RNF04 -- Autenticação e autorização:} o acesso deve ser controlado por autenticação \textit{stateless} baseada em tokens JWT, com permissões distintas para os perfis gestor e motorista.
  \item \textbf{RNF05 -- Privacidade:} o sistema deve tratar apenas os dados pessoais estritamente necessários à operação (nome, e-mail, telefone e, opcionalmente, número da CNH), sem dados financeiros ou sensíveis, em atenção ao princípio da necessidade previsto no art.~6º, inciso III, da LGPD.
  \item \textbf{RNF06 -- Conteinerização:} os ambientes de desenvolvimento e de implantação devem ser padronizados com Docker e Docker Compose.
  \item \textbf{RNF07 -- Desempenho:} as requisições à API devem ser respondidas em menos de 300~ms.
\end{itemize}

% ---------------------------------------------------------------------
% 3. Justificativa das escolhas tecnológicas (seção de arquitetura)
% ---------------------------------------------------------------------

A adoção de TypeScript tanto no \textit{front-end} quanto no \textit{back-end} permite que uma única linguagem seja utilizada em toda a aplicação, reduzindo o custo cognitivo de manutenção em um projeto desenvolvido individualmente. No \textit{back-end}, o NestJS foi escolhido por impor nativamente a separação em camadas e a injeção de dependências, o que torna a arquitetura \textit{Controller}--\textit{Service}--\textit{Repository} parte da estrutura do próprio \textit{framework}, e não apenas uma convenção do desenvolvedor.

Para o motorista, optou-se por uma aplicação web progressiva em vez de um aplicativo nativo. As funcionalidades exigidas em campo (consultar as coletas do dia e atualizar seu status) não dependem de recursos exclusivos de plataformas nativas, e o acesso por meio de um endereço web elimina a etapa de instalação por lojas de aplicativos. Essa característica é relevante tanto para a adoção quanto para a validação empírica, pois o sistema pode ser disponibilizado aos participantes pelo mesmo canal que se pretende substituir.

% ---------------------------------------------------------------------
% 4. Metodologia — ajuste no plano de validação
%    Use se o MVP for apresentado nas sessões de validação da solução.
% ---------------------------------------------------------------------

A validação ocorre em duas etapas. A validação do problema é realizada por meio de entrevistas semiestruturadas com três a cinco profissionais do setor e de um questionário quantitativo, com o objetivo de verificar a ocorrência e a frequência das dificuldades descritas. A validação da solução é realizada por meio de testes de usabilidade com o MVP, nos quais gestores e motoristas executam tarefas representativas sem treinamento prévio, sendo registradas a taxa de sucesso e o tempo de cada tarefa. Ao final de cada sessão, aplica-se o questionário \textit{System Usability Scale} (SUS), adotando-se como meta uma pontuação igual ou superior a 68, valor correspondente à média observada em estudos de usabilidade.
